% !TeX root = closed-systems-from-comparison-completeness.tex
% ============================================================
\chapter{Intrinsic Numerical Closure, Unified Scalar Channel, and Same--Channel Corrections}
\label{ch:numerical-closure}
\label{ch:numbers-shell-closure}
% ============================================================

% ============================================================
\section{Purpose and logical status}
\label{sec:shell-purpose}
% ============================================================

This chapter closes the intrinsic numerical layer of the closed relational
theory.
Its role is to separate what transport-closed structure fixes internally
from what remains a distinct realization choice.

No observer-side semantic input enters the core closure statements.
In particular, no external empirical normalization is required to state the
intrinsic scalar law itself.

The objectives are:
\begin{enumerate}[label=\textup{(\roman*)}]
\item classify intrinsic scalar observables carried by the stabilized
      quadratic sector;
\item identify the single normalization orbit of that scalar channel;
\item constrain admissible subleading scalar corrections;
\item separate intrinsic sector observables from downstream physical
      interpretation maps.
\end{enumerate}

The chapter proceeds in two layers.
First comes the forced intrinsic closure chain for scalar depth laws.
Second comes a compatibility bridge that preserves legacy shell labels
without introducing new observer-side data.

% ============================================================
\section{Structural inputs from Chapters~19--20}
\label{sec:shell-inputs}
% ============================================================

We use the stabilized quadratic carrier and depth architecture established in
\cref{ch:matter,ch:matter-numbers}:
\[
  \mathcal K\cong F^2/F^3,
  \qquad
  Q:\mathcal K_{\mathbb R}\to\mathbb R_{\ge0},
  \qquad
  k\in\mathbb N,
\]
together with the quartic depth polynomial
\[
  q(k):=4\binom{k}{4}
\]
from \cref{def:numbers-qk}.

We write
\[
  \Sigma(k):=Q(\kappa_k)
\]
for the intrinsic scalar profile indexed by depth.

Methodologically, each statement below is typed as either
\emph{intrinsic/forced} or \emph{compatibility/conditional}.
This mirrors the bookkeeping discipline used in
\cref{ch:matter-numbers}.

\begin{proposition}[Inherited dependency role in intrinsic form]
\label{prop:shell-dependency-role}
The unconditional intrinsic statements in this chapter use two saved
upstream inputs:
\begin{enumerate}[label=\textup{(\roman*)}]
\item the Chapter~19 scalar-channel framework
      \cref{thm:quadratic-scalar,thm:spectral-scalar-identification,thm:spectral-scalar-invariant};
\item the Chapter~20 quartic-backbone and rigidity-envelope package,
      read with its saved distinction between structural consequences
      and extra compatibility inputs.
\end{enumerate}
Exact identification of a chosen scalar profile with
\(q(k)=4\binom{k}{4}\), exact ratio formulas for that profile, and any
observer-side or Hilbert-branch realization statement require
additional compatibility or realization hypotheses beyond those
unconditional inputs.
\end{proposition}

\begin{proof}
Chapter~19 fixes the unique scalar channel up to positive
normalization, the canonical factorization of sector scalars through
that channel, and transport covariance.
Chapter~20 then supplies the quartic combinatorial backbone
\(q(k)=4\binom{k}{4}\) for the named quartic observable together with
the general quartic-growth and quartic-leading-law consequences used
for admissible scalar depth observables.
Those inputs support the unconditional intrinsic statements of this
chapter at exactly that strength.
Whenever the chapter speaks about exact equality with the quartic
backbone, exact ratio formulas for a chosen mass-coordinate
presentation, or the later Hilbert-branch/operator route, those claims
must be typed explicitly as compatibility-level or realization-level
statements rather than folded into the unconditional intrinsic
package.
\end{proof}

% ============================================================
\section{Intrinsic scalar closure chain}
\label{sec:shell-intrinsic-closure}
% ============================================================

This section begins with the forced intrinsic consequences of the
Chapter~19/20 package and then marks later compatibility or
conditional refinements explicitly where they enter.

\begin{proposition}[Scaling covariance]
\label{prop:scaling}
If \(\Sigma\) is admissible and \(\lambda>0\), then
\[
  \Sigma_\lambda(k):=\lambda\Sigma(k)
\]
is admissible.
\end{proposition}

\begin{proof}
Admissibility is stable under positive rescaling, since the scalar channel
is unique only up to positive normalization
(\cref{thm:quadratic-scalar,thm:spectral-ratio-properties}).
\end{proof}

\begin{theorem}[Intrinsic scalar-channel closure]
\label{thm:closure}
Every admissible intrinsic scalar profile along the depth architecture
remains in the unique scalar channel fixed by
\cref{thm:quadratic-scalar,thm:spectral-scalar-identification}.
Moreover, its leading growth is quartic in the sense that
\(\Sigma(k)=O(k^4)\), and under the Chapter~20 rigidity hypotheses its
universal leading law is quartic.
\end{theorem}

\begin{proof}
Channel uniqueness is exactly the Chapter~19 scalar-channel statement.
The quartic growth envelope and quartic leading law are the saved
Chapter~20 consequences recorded after patch P-218; they do not by
themselves identify the full profile with the exact quartic backbone.
\end{proof}

\begin{corollary}[Conditional ratio rigidity for the quartic backbone]
\label{cor:ratio-rigidity}
Assume the chosen intrinsic scalar profile is identified with the
Chapter~20 quartic observable up to one normalization constant, so that
\(\Sigma(k)=\lambda q(k)\) for some \(\lambda>0\).
Then for any admissible \(k_1,k_2\) with \(q(k_2)\neq0\),
\[
  \frac{\Sigma(k_1)}{\Sigma(k_2)}
  =
  \frac{q(k_1)}{q(k_2)}.
\]
\end{corollary}

\begin{proof}
Under the stated identification, the common factor \(\lambda\) cancels.
\end{proof}

\begin{theorem}[Uniqueness of scalar channel]
\label{thm:scalar-channel}
All intrinsic scalar observables descending from the stabilized quadratic
carrier factor through one and the same one-dimensional scalar channel.
\end{theorem}

\begin{proof}
This is precisely the channel-uniqueness content of
\cref{thm:quadratic-scalar,thm:spectral-scalar-identification}.
\end{proof}

\begin{theorem}[No absolute normalization]
\label{thm:no-absolute}
No intrinsic construction in Chapters~1--21 fixes the overall
normalization constant \(\lambda\).
\end{theorem}

\begin{proof}
By \cref{thm:spectral-ratio-properties}, intrinsic structure fixes ratios,
whereas absolute scale requires additional realization data.
\end{proof}

% ============================================================
\section{Correction structure inside the same channel}
\label{sec:shell-corrections}
% ============================================================

With the scalar channel fixed, we now record the compatibility model
used later when one discusses subleading depth dependence without
creating a new invariant type.

\begin{theorem}[No new scalar invariants]
\label{thm:no-new}
Admissible scalar corrections do not create an independent scalar channel;
they remain same-channel deformations of the intrinsic quartic profile.
\end{theorem}

\begin{proof}
By \cref{thm:scalar-channel}, the scalar carrier is one-dimensional.
Hence admissible scalar corrections may change depth dependence, but not the
underlying scalar-channel type.
\end{proof}

\begin{theorem}[Conditional binomial correction model]
\label{thm:basis}
Assume a same-channel scalar profile \(M(k)\) is written as
\[
  M(k)=\lambda\,4\binom{k}{4}+R_{\le2}(k),
\]
where \(\lambda\in\mathbb R\) and the remainder \(R_{\le2}(k)\) is a discrete
polynomial of degree at most \(2\).
Then there exist constants \(\alpha,\beta,\gamma,\delta\in\mathbb R\)
such that
\[
  M(k)
  =
  \alpha\binom{k}{4}
  +
  \beta\binom{k}{2}
  +
  \gamma k
  +
  \delta.
\]
\end{theorem}

\begin{proof}
Because \(R_{\le2}(k)\) has degree at most \(2\), it can be written in the
binomial basis \(\beta\binom{k}{2}+\gamma k+\delta\).
Absorbing the quartic backbone coefficient into
\(\alpha:=4\lambda\) yields the displayed form.
\end{proof}

\begin{proposition}[Conditional leading structure]
\label{prop:leading-structure}
Under the hypotheses of \cref{thm:basis},
\[
  M(k)=\alpha\binom{k}{4}+O(k^2)
  \qquad (k\to\infty).
\]
In particular, the quartic backbone controls the leading asymptotic
behavior.
\end{proposition}

\begin{proof}
The remainder terms \(\beta\binom{k}{2}+\gamma k+\delta\) have degree
at most \(2\), so they contribute only \(O(k^2)\).
\end{proof}

\begin{theorem}[Conditional cubic rigidity under the correction model]
\label{thm:cubic}
Under the hypotheses of \cref{thm:basis}, the cubic coefficient of
\(M(k)\) is fixed entirely by the quartic backbone and cannot be
independently deformed by the modeled remainder.
\end{theorem}

\begin{proof}
The quartic backbone \(4\binom{k}{4}\) already determines its own cubic
term. Since the modeled remainder has degree at most \(2\), it cannot
change that cubic coefficient.
\end{proof}

\begin{proposition}[Conditional structure of modeled corrections]
\label{prop:correction-structure}
Under the hypotheses of \cref{thm:basis}, modeled same-channel
corrections first appear at quadratic order and below.
\end{proposition}

\begin{proof}
This is exactly the degree statement encoded in the assumption on
\(R_{\le2}(k)\) and restated by \cref{thm:cubic}.
\end{proof}

% ============================================================
\section{Intrinsic spectral mass law, ratio rigidity, and asymptotic exactness}
\label{sec:intrinsic-spectral-mass-law}
% ============================================================

This section rewrites the intrinsic scalar closure chain in spectral-mass
language and isolates the asymptotic structure of realization-level
corrections.
At intrinsic level, no new input is added beyond the scalar-channel and
quartic-depth package
\cref{thm:quadratic-scalar,thm:spectral-scalar-identification,thm:spectral-ratio-properties,thm:numbers-quartic-formula,def:numbers-qk}.
Its statements are explicitly split into two types:
\begin{enumerate}[label=\textup{(\roman*)}]
\item \emph{intrinsic/forced}: the one-channel structural facts,
      quartic leading behavior, and absence of intrinsic absolute
      normalization already fixed by
      \cref{thm:closure,thm:no-absolute};
\item \emph{compatibility/conditional}: exact quartic-backbone formulas
      for a chosen mass-coordinate presentation, together with any
      consequences that require additional
      Hilbert-branch realization input beyond Chapters~1--21, namely the
      operator-route package
      \cref{thm:intrinsic-spectral-mass-law,thm:exact-ratio-law,thm:hilbert-realization-sector-scalar,lem:sector-quadratic-form,prop:matter-realization-form-route}.
\end{enumerate}

Define the quartic backbone
\[
  B_4(k):=4\binom{k}{4}
  =
  \frac{k(k-1)(k-2)(k-3)}{6},
  \qquad
  k\in\mathbb N.
\]
By \cref{def:numbers-qk}, this is exactly the intrinsic depth law
\(q(k)\).

\begin{definition}[Intrinsic mass-coordinate profile]
\label{def:intrinsic-mass-coordinate-profile}
Write
\[
  m_{\mathrm{int}}(k):=\Sigma(k).
\]
This is an intrinsic scalar profile (mass-coordinate language), not yet an
observer-side measured mass assignment.
Equivalently, it is the depth-indexed presentation of the same canonical
scalar channel fixed by
\cref{thm:quadratic-scalar,thm:spectral-scalar-identification}.
\end{definition}

\begin{theorem}[Conditional intrinsic spectral mass law]
\label{thm:intrinsic-spectral-mass-law}
Assume the mass-coordinate profile is the Chapter~20 quartic backbone
written in scalar-channel units, namely
\[
  m_{\mathrm{int}}(k)=\lambda\,4\binom{k}{4}
\]
for some \(\lambda>0\).
Then the displayed quartic formulas, exact table, and exact ratio laws
below hold for that chosen compatibility model.
\end{theorem}

\begin{proof}
This is a direct restatement of the assumed identification with the
exact quartic backbone \(q(k)=4\binom{k}{4}\).
\end{proof}

\begin{proposition}[Conditional unnormalized intrinsic mass table]
\label{prop:unnormalized-intrinsic-mass-table}
Under \cref{thm:intrinsic-spectral-mass-law}, the quartic backbone and
mass-coordinate profile are given by the displayed table for
\(k=0,\dots,20\):
\[
\begin{array}{c|c|c}
k & B_4(k) & m_{\mathrm{int}}(k) \\
\hline
0 & 0 & 0 \\
1 & 0 & 0 \\
2 & 0 & 0 \\
3 & 0 & 0 \\
4 & 4 & 4\lambda \\
5 & 20 & 20\lambda \\
6 & 60 & 60\lambda \\
7 & 140 & 140\lambda \\
8 & 280 & 280\lambda \\
9 & 504 & 504\lambda \\
10 & 840 & 840\lambda \\
11 & 1320 & 1320\lambda \\
12 & 1980 & 1980\lambda \\
13 & 2860 & 2860\lambda \\
14 & 4004 & 4004\lambda \\
15 & 5460 & 5460\lambda \\
16 & 7280 & 7280\lambda \\
17 & 9520 & 9520\lambda \\
18 & 12240 & 12240\lambda \\
19 & 15504 & 15504\lambda \\
20 & 19380 & 19380\lambda
\end{array}
\]
\end{proposition}

\begin{proof}
Direct evaluation of
\(B_4(k)=4\binom{k}{4}\), followed by
\(m_{\mathrm{int}}(k)=\lambda B_4(k)\)
from \cref{thm:intrinsic-spectral-mass-law}.
\end{proof}

\begin{theorem}[Conditional exact ratio law]
\label{thm:exact-ratio-law}
Under \cref{thm:intrinsic-spectral-mass-law}, for any depths
\(k,\ell\) with \(B_4(\ell)\neq0\),
\[
  \frac{m_{\mathrm{int}}(k)}{m_{\mathrm{int}}(\ell)}
  =
  \frac{B_4(k)}{B_4(\ell)}
  =
  \frac{\binom{k}{4}}{\binom{\ell}{4}}.
\]
\end{theorem}

\begin{proof}
Immediate from \cref{thm:intrinsic-spectral-mass-law}.
\end{proof}

\begin{corollary}[Conditional sample exact ratio predictions]
\label{cor:sample-exact-ratio-predictions}
Under \cref{thm:intrinsic-spectral-mass-law}, the displayed sample
quartic-backbone ratios follow:
\[
  \frac{m_{\mathrm{int}}(5)}{m_{\mathrm{int}}(4)}=5,
  \qquad
  \frac{m_{\mathrm{int}}(6)}{m_{\mathrm{int}}(4)}=15,
  \qquad
  \frac{m_{\mathrm{int}}(18)}{m_{\mathrm{int}}(6)}=204,
  \qquad
  \frac{m_{\mathrm{int}}(35)}{m_{\mathrm{int}}(6)}=\frac{10472}{3}.
\]
\end{corollary}

\begin{proof}
Substitute into
\cref{thm:exact-ratio-law}.
\end{proof}

\begin{theorem}[Normalization obstruction]
\label{thm:normalization-obstruction-final}
The closed relational system does not determine the absolute normalization
scale of \(m_{\mathrm{int}}\).
\end{theorem}

\begin{proof}
This is exactly
\cref{thm:no-absolute}:
intrinsic structure fixes only dimensionless ratios
(\cref{thm:spectral-ratio-properties}) and leaves absolute normalization to
separate realization data
(\cref{rem:numbers-forced-vs-calibrated,sec:cond-open}).
\end{proof}

\begin{corollary}[Intrinsic prediction content at saved strength]
\label{cor:complete-intrinsic-prediction-content}
At unconditional intrinsic level, the theory determines the unique
scalar channel, quartic leading behavior, and the absence of intrinsic
absolute normalization.
If one additionally identifies the mass-coordinate profile with the
Chapter~20 quartic backbone up to normalization, then the exact
quartic formulas and exact ratio consequences above follow.
\end{corollary}

\begin{proof}
The unconditional intrinsic claims are exactly the one-channel and
normalization statements from
\cref{thm:closure,thm:no-absolute}.
The conditional quartic-backbone claims are
\cref{thm:intrinsic-spectral-mass-law,thm:exact-ratio-law}.
\end{proof}

\begin{remark}[Operator status and realization boundary]
\label{rem:intrinsic-spectral-law-operator-status}
The operator statement needed to realize
\(m_0(\sigma)=Q(\kappa_\sigma)\) as Hilbert-branch spectral data is the
target operator-route package
\cref{thm:hilbert-realization-sector-scalar,lem:sector-quadratic-form,prop:matter-realization-form-route}.
Within Chapters~1--21, this remains a realization-level step rather than an
already completed intrinsic theorem.
\end{remark}

\begin{definition}[Depth operator \textup{(conditional realization datum)}]
\label{def:depth-operator}
Assume a realized Hilbert branch \(H_{\mathrm{quad}}\) with a depth-indexed
sector family \(\{\psi_k\}_{k\in\mathbb N}\subset H_{\mathrm{quad}}\), each
\(\psi_k\neq 0\), compatible with the intrinsic depth label
\(k\) from \cref{thm:spectral-depth-separation}.
Define the depth operator \(\mathsf K\) on the finite span
\(D_{\mathrm{fin}}:=\operatorname{span}\{\psi_k\}\) by
\[
  \mathsf K\psi_k=k\,\psi_k.
\]
\end{definition}

\begin{definition}[Quartic backbone operator \textup{(conditional)}]
\label{def:quartic-backbone-operator}
On \(D_{\mathrm{fin}}\), define
\[
  B_4(\mathsf K)
  :=
  4\binom{\mathsf K}{4}
  =
  \frac{\mathsf K(\mathsf K-1)(\mathsf K-2)(\mathsf K-3)}{6}.
\]
\end{definition}

\begin{definition}[Realization defect operator \textup{(conditional model)}]
\label{def:realization-defect-operator-final}
Assume a realized sector operator \(\mathsf S_{\mathrm{real}}\) on
\(D_{\mathrm{fin}}\) admits
the decomposition
\[
  \mathsf S_{\mathrm{real}}=\lambda B_4(\mathsf K)+R_{\mathrm{def}},
\]
where \(R_{\mathrm{def}}\) is the realization defect operator.
\end{definition}

\begin{proposition}[Realized mass decomposition]
\label{prop:realized-mass-decomposition}
Under
\cref{def:realization-defect-operator-final},
for every depth vector \(\psi_k\in D_{\mathrm{fin}}\),
\[
  m_{\mathrm{real}}(k)
  :=
  \langle \psi_k,\mathsf S_{\mathrm{real}}\psi_k\rangle
  =
  \lambda B_4(k)+\langle \psi_k,R_{\mathrm{def}}\psi_k\rangle.
\]
\end{proposition}

\begin{proof}
Insert
\(\mathsf S_{\mathrm{real}}=\lambda B_4(\mathsf K)+R_{\mathrm{def}}\)
and use
\(B_4(\mathsf K)\psi_k=B_4(k)\psi_k\).
\end{proof}

\begin{definition}[Depth renormalization parameter]
\label{def:depth-renormalization-parameter-final}
Assume \(B_4'(k)\neq 0\).
Define
\[
  \delta_k
  :=
  \frac{\langle \psi_k,R_{\mathrm{def}}\psi_k\rangle}{\lambda B_4'(k)}.
\]
\end{definition}

\begin{proposition}[First-order corrected spectrum]
\label{prop:first-order-corrected-spectrum}
Under
\cref{def:depth-renormalization-parameter-final},
\[
  m_{\mathrm{real}}(k)
  =
  \lambda\bigl(B_4(k)+B_4'(k)\delta_k\bigr),
\]
and hence
\[
  m_{\mathrm{real}}(k)
  =
  \lambda B_4(k+\delta_k)
  +
  O\!\bigl(\delta_k^2B_4''(k)\bigr)
\]
to first order in realization defect.
\end{proposition}

\begin{proof}
By
\cref{def:depth-renormalization-parameter-final},
\(
\langle \psi_k,R_{\mathrm{def}}\psi_k\rangle=\lambda B_4'(k)\delta_k
\).
Substitute into
\cref{prop:realized-mass-decomposition}.
Then apply Taylor expansion of the quartic polynomial \(B_4\).
\end{proof}

\begin{lemma}[Quartic derivative asymptotics]
\label{lem:quartic-derivative-asymptotics-final}
For
\[
  B_4(k)
  =
  4\binom{k}{4}
  =
  \frac{k(k-1)(k-2)(k-3)}{6},
\]
one has
\[
  B_4'(k)
  =
  \frac{4k^3-18k^2+22k-6}{6}
  =
  \frac23\,k^3+O(k^2),
\]
and
\[
  B_4''(k)
  =
  \frac{12k^2-36k+22}{6}
  =
  2k^2+O(k).
\]
\end{lemma}

\begin{proof}
Differentiate the polynomial expression for \(B_4(k)\) directly.
\end{proof}

\begin{theorem}[Asymptotic exactness of the quartic mass law]
\label{thm:asymptotic-exactness-quartic-mass-law-final}
Assume
\[
  \langle \psi_k,R_{\mathrm{def}}\psi_k\rangle = O(k^p)
  \qquad
  \text{for some }p<3.
\]
Then:
\begin{enumerate}[label=\textup{(\roman*)}]
\item
      \[
        \delta_k = O(k^{p-3}),
        \qquad
        \delta_k\to 0
        \quad (k\to\infty);
      \]
\item
      \[
        m_{\mathrm{real}}(k)=\lambda B_4(k)+o(B_4(k));
      \]
\item
      \[
        \frac{m_{\mathrm{real}}(k)}{\lambda B_4(k)}\to 1
        \qquad
        (k\to\infty).
      \]
\end{enumerate}
In particular, if
\(
\langle \psi_k,R_{\mathrm{def}}\psi_k\rangle = O(k^2)
\), then
\(
\delta_k = O(k^{-1})
\).
\end{theorem}

\begin{proof}
By
\cref{def:depth-renormalization-parameter-final},
\[
  \delta_k
  =
  \frac{\langle \psi_k,R_{\mathrm{def}}\psi_k\rangle}{\lambda B_4'(k)}.
\]
By hypothesis,
\(
\langle \psi_k,R_{\mathrm{def}}\psi_k\rangle = O(k^p)
\), and by
\cref{lem:quartic-derivative-asymptotics-final},
\(
B_4'(k)=\Theta(k^3)
\).
Hence
\(
\delta_k = O(k^{p-3})
\), so \(\delta_k\to0\) because \(p<3\).

Next,
\[
  m_{\mathrm{real}}(k)
  =
  \lambda B_4(k)+\langle \psi_k,R_{\mathrm{def}}\psi_k\rangle
  =
  \lambda B_4(k)+O(k^p).
\]
Since
\(
B_4(k)=\frac16k^4+O(k^3)=\Theta(k^4)
\)
and \(p<3<4\), one has
\(
O(k^p)=o(B_4(k))
\).
Therefore
\(
m_{\mathrm{real}}(k)=\lambda B_4(k)+o(B_4(k))
\), and dividing by \(\lambda B_4(k)\) gives item \textup{(iii)}.
The final sentence is the case \(p=2\).
\end{proof}

\begin{proposition}[Conditional exact-ratio test for a chosen depth assignment]
\label{prop:exact-ratio-test}
Assume the chosen mass-coordinate profile is identified with the
quartic backbone as in \cref{thm:intrinsic-spectral-mass-law}.
For any proposed depth assignment
\[
  k_1,\dots,k_n,
\]
that compatibility model then predicts exact ratios
\[
  \frac{m_{\mathrm{int}}(k_i)}{m_{\mathrm{int}}(k_j)}
  =
  \frac{B_4(k_i)}{B_4(k_j)}.
\]
Hence the quartic-backbone compatibility model admits an exact
unnormalized ratio test.
\end{proposition}

\begin{proof}
Immediate from \cref{thm:exact-ratio-law}.
\end{proof}

\begin{remark}[Interpretive consequence]
\label{rem:exact-ratio-interpretive-consequence}
Under the same quartic-backbone identification, if observed realized
ratios for a chosen depth assignment are close to the exact quartic
ratios, this supports using that quartic backbone as the leading
intrinsic model for the chosen assignment.
Residual discrepancy is then naturally modeled as lower-order realization
defect, and under
\cref{thm:asymptotic-exactness-quartic-mass-law-final} this defect is
asymptotically negligible.
\end{remark}

\begin{theorem}[Complete intrinsic spectral prediction at the saved split]
\label{thm:complete-intrinsic-spectral-prediction}
The closed relational system determines the scalar channel and its
quartic leading behavior, but not an intrinsic absolute normalization.
Under the explicit quartic-backbone identification used in this
section, one further obtains the exact table and exact ratio formulas
for that compatibility model.
Under an additional subcubic realization-defect bound, the realized
spectrum is asymptotically exact relative to that quartic backbone in
the sense of
\cref{thm:asymptotic-exactness-quartic-mass-law-final}.
\end{theorem}

\begin{proof}
The unconditional intrinsic statement is exactly the one-channel and
quartic-leading-behavior package summarized in
\cref{thm:closure,thm:no-absolute}.
The conditional quartic-backbone table and ratio formulas are
\cref{thm:intrinsic-spectral-mass-law,prop:unnormalized-intrinsic-mass-table,thm:exact-ratio-law}.
The final statement is
\cref{thm:asymptotic-exactness-quartic-mass-law-final}.
\end{proof}

% ============================================================
\section{Handoff to matter-sector realization}
\label{sec:shell-handoff}
% ============================================================

This chapter closes intrinsic scalar closure at the level of quartic
leading behavior and a single normalization orbit, while retaining
same-channel correction structure only as a compatibility bridge for
modeled same-channel corrections.
The next chapter (
\cref{ch:matter-realization}
) handles branch-level realization of the canonical sector invariant
\(m_0(\sigma)=Q(\kappa_\sigma)\).
Its first remaining theorem target is the Hilbert-branch operator
realization
\cref{thm:hilbert-realization-sector-scalar}, via the closed positive
form route.
It then isolates the explicit boundary between intrinsic sector observables
and observer-side interpretation maps.


% ============================================================
\subsection{Forced-upgrade roadmap for the operator target}
\label{subsec:shell-forced-upgrade-roadmap}
% ============================================================

A full promotion of the Chapter~22 operator target
\cref{thm:hilbert-realization-sector-scalar}
from realization/conditional status to intrinsic/forced status requires a
theorem chain proved from existing theorem-level data only.
The following labeled items record the minimal roadmap targets whose later
discharge would suffice for that upgrade: canonical sector-to-Hilbert
embedding, induced sector form construction, and closed-form closability.

\begin{theorem}[Roadmap target: canonical sector-to-Hilbert embedding]
\label{thm:roadmap-sector-to-hilbert-embedding}
To promote the Chapter~22 operator statement to intrinsic/forced status,
it would suffice to prove the following statement from only
\cref{thm:interface-stabilized,ch:hilbert,def:spectral-defect-class,def:spectral-scalar,thm:spectral-depth-separation}:
there exists a canonical linear map
\[
  \iota_{\mathrm{sec}}:
  \operatorname{span}\{\kappa_\sigma:\sigma\text{ primitive excitation sector}\}
  \hookrightarrow
  H_{\mathrm{quad}}
\]
with
\(\iota_{\mathrm{sec}}(\kappa_\sigma)=\psi_\sigma\neq0\),
compatible with intrinsic depth labels and transport covariance.
\end{theorem}

\begin{lemma}[Roadmap target: embedded sector form realizes intrinsic scalar values]
\label{lem:roadmap-embedded-sector-form}
Assuming \cref{thm:roadmap-sector-to-hilbert-embedding}, the later
intrinsic upgrade would also be secured by proving the following target
statement from
\cref{thm:quadratic-scalar,thm:spectral-scalar-identification,prop:hilbert-real-inner}:
define on the finite sector span
\[
  q_{\mathrm{sec}}(\psi,\phi)
  :=
  \left\langle
    \iota_{\mathrm{sec}}^{-1}\psi,
    \iota_{\mathrm{sec}}^{-1}\phi
  \right\rangle_Q.
\]
Then \(q_{\mathrm{sec}}\) is positive and densely defined, and for each
primitive sector \(\sigma\),
\[
  q_{\mathrm{sec}}(\psi_\sigma,\psi_\sigma)
  =
  m_0(\sigma)\,\|\psi_\sigma\|^2.
\]
Hence \(q_{\mathrm{sec}}\) would be of the input type required by
\cref{lem:sector-quadratic-form}.
\end{lemma}

\begin{theorem}[Roadmap target: closed-form closability from intrinsic control]
\label{thm:roadmap-closed-form-closability}
Assuming \cref{lem:roadmap-embedded-sector-form}, the later intrinsic
upgrade would further be secured by proving this target statement:
if intrinsic Hilbert-branch control estimates imply that
\(q_{\mathrm{sec}}\) is closable, then its closure
\(\overline q_{\mathrm{sec}}\) is a closed positive form on
\(H_{\mathrm{quad}}\).
This is exactly the closed/closable-form hypothesis needed in
\cref{prop:matter-realization-form-route}.
\end{theorem}

\begin{proposition}[Roadmap implication for a later forced upgrade of the Chapter~22 operator target]
\label{prop:roadmap-implies-ch22-operator-target}
If
\cref{thm:roadmap-sector-to-hilbert-embedding,lem:roadmap-embedded-sector-form,thm:roadmap-closed-form-closability}
were later discharged from existing theorem-level inputs alone, then
that discharge would supply the sufficiency route through
\cref{prop:matter-realization-form-route} for promoting
\cref{thm:hilbert-realization-sector-scalar} to intrinsic/forced status.
\end{proposition}

\begin{proof}
A later proof of the first two roadmap targets would provide the
sector-form data required by \cref{lem:sector-quadratic-form}.
A later proof of the third would provide the closed/closable-form
hypothesis used by \cref{prop:matter-realization-form-route}.
Together, those later discharges would therefore yield the canonical
positive self-adjoint operator with sector eigenvalue law
\(\mathsf S\psi_\sigma=m_0(\sigma)\psi_\sigma\), which is exactly the
content that would promote
\cref{thm:hilbert-realization-sector-scalar} to intrinsic/forced
status.
\end{proof}

\begin{corollary}[Forced-upgrade criterion recorded at roadmap strength]
\label{cor:roadmap-forced-upgrade-criterion}
A full promotion from realization/conditional to intrinsic/forced for the
Chapter~22 operator statement is justified exactly when the roadmap
targets
\cref{thm:roadmap-sector-to-hilbert-embedding,lem:roadmap-embedded-sector-form,thm:roadmap-closed-form-closability}
are later proved without adding new observer-side assumptions.
\end{corollary}

\begin{proof}
Immediate from
\cref{prop:roadmap-implies-ch22-operator-target,sec:cond-open}.
\end{proof}

% ============================================================
\section{Compatibility bridge for legacy shell labels}
\label{sec:shell-legacy-compat}
% ============================================================

For consistency with the existing ledger and cross-chapter references, we
record a compact compatibility layer in legacy shell notation.
These labels are optional interface markers and do not add observer-side
inputs to the intrinsic core established above.
They function as interface-level translators between the current closure
language and older shell bookkeeping.

\begin{definition}[Signed shell profile (compatibility form)]
\label{def:shell-signed-profile}
At fixed depth, a signed shell profile is a finite family
\((\Delta_\eta(\sigma))_{\eta\in\Lambda_\sigma}\) with
\(\Delta_\eta(\sigma)\in\mathbb Z\).
\end{definition}

\begin{definition}[Shell splitting invariant (compatibility form)]
\label{def:shell-S}
Define
\[
  S(\sigma):=\sum_{\eta\in\Lambda_\sigma}\Delta_\eta(\sigma)^2.
\]
\end{definition}

\begin{definition}[Shell completeness property (legacy label)]
\label{def:shell-completeness}
A realized branch is \emph{shell-complete} at fixed depth when
shell-visible scalar distinction factors through \(S(\sigma)\).
\end{definition}

\begin{theorem}[Signed-permutation criterion for shell completeness]
\label{thm:shell-completeness-from-symmetry}
If the shell residual at depth \(k\) is quadratic and invariant under
independent sign flips and coordinate permutations of shell coordinates,
then it is proportional to
\(\sum_\eta\Delta_\eta^2\).
If the proportionality coefficient is depth-independent, then
the shell-completeness property of \cref{def:shell-completeness} holds.
\end{theorem}

\begin{proof}
The invariance forces diagonal isotropy of the associated symmetric bilinear
form, so the quadratic residual is a scalar multiple of the Euclidean square
sum.
Depth-independent proportionality then yields a single universal shell
channel.
\end{proof}

\begin{corollary}[Conditional discharge of shell completeness]
\label{cor:shell-completeness-discharged}
Whenever \cref{thm:shell-completeness-from-symmetry} holds with
depth-independent coefficient, the shell-completeness property
\cref{def:shell-completeness} need not be stipulated independently.
\end{corollary}

\begin{proof}
Immediate from \cref{thm:shell-completeness-from-symmetry}.
\end{proof}

\begin{theorem}[Transport-to-shell derivation of symmetry hypotheses]
\label{thm:shell-symmetry-from-transport}
Fix depth \(k\) and an interface map
\(I_k:\mathbb R^{\Lambda_k}\to\mathcal K_{\mathbb R}\).
If
\[
  R_k(\Delta)=Q(I_k\Delta)
\]
and for each signed permutation \(g\) there is a realized transport action
\(T_g\) on \(\mathcal K_{\mathbb R}\) with
\[
  I_k(g\Delta)=T_gI_k(\Delta),
  \qquad
  Q(T_gx)=Q(x),
\]
then \(R_k\) satisfies the symmetry hypotheses of
\cref{thm:shell-completeness-from-symmetry}.
\end{theorem}

\begin{proof}
Linearity of \(I_k\) and quadraticity of \(Q\) make \(R_k\) quadratic.
Intertwining together with \(Q\)-invariance implies invariance of
\(R_k\) under the signed-permutation action.
Hence the symmetry hypotheses are satisfied.
\end{proof}

\begin{theorem}[Interface checklist for transport-derived symmetry]
\label{thm:shell-interface-checklist}
To verify condition \textup{(iii)} of
\cref{thm:shell-symmetry-from-transport} in-model, it suffices to check
intertwining and \(Q\)-invariance on finite generators:
coordinate sign flips \(s_i\) and adjacent transpositions \(\tau_i\),
together with their signed-permutation relations.
\end{theorem}

\begin{proof}
These generators and relations present the hyperoctahedral group.
If intertwining and \(Q\)-invariance hold on generators, they extend to all
group elements by composition.
\end{proof}

\begin{corollary}[Transport-derived shell-completeness route]
\label{cor:shell-completeness-from-transport}
If \cref{thm:shell-symmetry-from-transport} holds for each depth and the
induced proportionality coefficients are depth-independent, then
the shell-completeness property \cref{def:shell-completeness} and
\cref{cor:shell-completeness-discharged} hold.
\end{corollary}

\begin{proof}
Combine
\cref{thm:shell-symmetry-from-transport,thm:shell-completeness-from-symmetry,cor:shell-completeness-discharged}.
\end{proof}

\begin{definition}[One-reference defect-evaluation datum (legacy label)]
\label{def:shell-defect-intrinsic}
A one-reference defect-evaluation datum is a non-balanced reference sector
with intrinsic defect data sufficient to recover the dimensionless shell
ratio parameter in one-reference form.
\end{definition}

\begin{definition}[One-reference anchor datum (legacy label)]
\label{def:shell-anchor}
A one-reference anchor datum is a reference sector with fixed intrinsic
scalar value chosen for absolute normalization in one-reference form.
\end{definition}

\begin{theorem}[Two-reference algebraic closure]
\label{thm:shell-two-reference-closure}
Assume a shell-corrected normalized law
\[
  m(\sigma)=\alpha\bigl(q(k_\sigma)+\gamma S(\sigma)\bigr)
\]
in a realized branch.
For two reference sectors \(\sigma_1,\sigma_2\) with known
\(m_i,k_i,S_i\) \((i=1,2)\), if
\(
  m_1S_2-m_2S_1\neq0
\), then
\[
  \gamma
  =
  \frac{m_2q(k_1)-m_1q(k_2)}{m_1S_2-m_2S_1},
\]
and
\[
  \alpha
  =
  \frac{m_1}{q(k_1)+\gamma S_1}
  =
  \frac{m_2}{q(k_2)+\gamma S_2}.
\]
\end{theorem}

\begin{proof}
Write the normalized law at \(\sigma_1,\sigma_2\), eliminate \(\alpha\) by
cross-multiplication, solve for \(\gamma\), and then substitute back to
recover \(\alpha\).
\end{proof}

\begin{corollary}[Discharge of one-reference closure data]
\label{cor:shell-two-reference-discharges}
Under \cref{thm:shell-two-reference-closure}, the one-reference data
\cref{def:shell-defect-intrinsic,def:shell-anchor} are optional.
\end{corollary}

\begin{proof}
Two-reference algebraic recovery determines both the dimensionless parameter
and the absolute normalization directly.
\end{proof}

% ============================================================
\section{Final closure statement}
\label{sec:shell-summary}
% ============================================================

\begin{theorem}[Maximal intrinsic closure]
\label{thm:final}
The closed system determines:
\begin{enumerate}[label=\textup{(\roman*)}]
\item quartic structural depth law up to one normalization orbit;
\item ratio rigidity and absence of intrinsic absolute normalization.
\end{enumerate}
In addition, under the chapter's explicit same-channel correction
model, the cubic coefficient is fixed by the quartic backbone and the
modeled remainder begins at quadratic order.
No additional intrinsic scalar degree of freedom is left implicit.
\end{theorem}

\begin{proof}
The intrinsic claims follow from
\cref{thm:closure,thm:spectral-ratio-properties,thm:no-absolute}.
The modeled correction statement follows from
\cref{thm:no-new,thm:basis,thm:cubic,prop:correction-structure}
under the explicit ansatz used in
\cref{sec:shell-corrections}.
\end{proof}

% ============================================================
\section{Conclusion}
\label{sec:shell-conclusion}
% ============================================================

The intrinsic numerical sector is now closed at theorem level:
the scalar law is fixed up to one normalization orbit, and any later
same-channel correction discussion remains an explicit compatibility
model rather than a new intrinsic scalar channel.

Accordingly, this chapter closes the intrinsic one-channel numerical
core of the main manuscript at saved strength: quartic leading
behavior and normalization-orbit freedom are fixed, while
same-channel correction models and realization-level identifications
remain explicit rather than forced.
The branch-level realization role of the canonical matter-sector
scalar is taken up next in \cref{ch:matter-realization}.